\section{Proofs for Section~\ref{sec:tails-compression}: Compressing the Pieces}
\label{app:compression}

We keep the conventions of \cref{app:tails}; in particular $k\ge1$ and
$N\ge2$. As in \cref{def:hole}, a \emph{comparison} of a tuple is a pair of
cuts of a host word and a target word such that the rest of the target is a
prefix of the rest of the host. A word
is \emph{primitive} if it is not a power of a shorter word, and two words
are \emph{conjugate} if one is a rotation of the other.

\subsection{Pieces}\label{app:compression-pieces}

Let $W$ be a tuple. A \emph{piece} of period $d$ is an interval $[s,e)$ of
positions of a word $w_l$ with $e-s\ge2d$ on which $w_l$ has period $d$,
whose first $d$ letters form a primitive word, its \emph{root}, and which is
maximal: $s=0$ or $w_l[s-1]\ne w_l[s-1+d]$, and $e=\len{w_l}$ or
$w_l[e]\ne w_l[e-d]$. Every factor of length at least $2p$ with period $p$
lies in a piece whose period is the least period of the factor, a divisor of
$p$.
\begin{enumerate}[label=\textup{(P\arabic*)}]
\item\label{it:P1} Two different pieces of period $d$ in one word overlap in
  fewer than $d$ positions, and their ends are more than $d$ apart: an
  overlap of $d$ positions would make the union an interval of period $d$.
\item\label{it:P2} Two pieces of one word with periods $d'\le d$ and roots
  that are not conjugate overlap in fewer than $d+d'$ positions: otherwise,
  by Fine and Wilf, the overlap has the period $\gcd(d,d')$, which is $d$
  and $d'$ because the roots are primitive, and the roots are conjugate.
\item\label{it:P3} Let $\alpha_d=Nd$ and $m_d=d\cdot N!$. Inside a piece
  $[s,e)$ of period $d$, $q_l(s+o+m_d)=q_l(s+o)$ whenever $\alpha_d\le o$ and
  $o+m_d\le e-s$. Indeed, reading the $d$ letters at the offsets
  $o',\dots,o'+d-1$ is the same map on states for all offsets $o'\equiv o$
  modulo $d$; its $N$th iterate lands on a cycle, whose length divides $N!$.
\end{enumerate}
Fix a primitive word $\zeta$, $d=\len\zeta$, and write
$\alpha,m,\beta,\kappa,\phi$ for $\alpha_d,m_d,\beta_d,\kappa_d,\phi_d$; so
$\beta=\alpha+m\ge4d$, $\kappa=2\beta+m$ and $\phi=6\beta+2m$. Let
$\mathcal C$ be the class of $\zeta$. For $P\in\mathcal C$ we write
$s_P,e_P$ for its ends and $L_P=e_P-s_P$.

\subsection{Admissible Deletions}\label{app:compression-admissible}

Let $(x_P)_{P\in\mathcal C}$ be an admissible deletion
(\cref{def:admissible}). The \emph{site} of $P$ is the interval
$[s_P+\beta,\,s_P+\beta+m\,x_P)$, and $U=(u_1,\dots,u_k)$ is the tuple
obtained by deleting all sites. We use the notation $\cmap_l$, $\gamma$ and
the facts on the deletion of positions from \cref{app:fresh-deletion},
writing $\cmap_l$ for the cut map of the word $w_l$; a \emph{junction} is
the image of a nonempty site. Put $L'_P=L_P-m\,x_P$.
\begin{enumerate}[label=\textup{(A\arabic*)}]
\item\label{it:A1} A nonempty site lies inside its piece, at least
  $\beta$ positions after its start and at least $5\beta+2m$ before its end.
  By \ref{it:P1}, sites of different pieces of one word are disjoint and not
  adjacent, and a position of a site lies in no other piece of the class.
\item\label{it:A2} By \ref{it:P3}, the two ends of a site carry the same
  state. So the state of $u_l$ at $\cmap_l(p)$ is the state of $w_l$ at $p$
  for every retained cut $p$; endpoints are preserved and no new state is
  visited.
\item\label{it:A3} For $P\in\mathcal C$ in the word $w_l$, the interval
  $P'=[\cmap_l(s_P),\cmap_l(e_P))$ is a piece of $U$ of the class
  $\mathcal C$, with the root of $P$ and the length $L'_P$; at an offset
  $o\le L'_P$ it carries the state that $P$ carries at the offset $o$, and
  at $o<L'_P$ the same letter. The map $P\mapsto P'$ is a bijection from
  $\mathcal C$ onto the pieces of $U$ of the class of $\zeta$.
\item\label{it:A4} If $o\le L_P$ and $o'\le L'_P$ are offsets with $o=o'$,
  or with $o,o'\ge\alpha$ and $o\equiv o'$ modulo $m$, then $P$ at $o$ and
  $P'$ at $o'$ carry the same state, by \ref{it:A3} and \ref{it:P3}.
\end{enumerate}
For \ref{it:A3}: the deleted block has a length divisible by $d$, so $P'$
has period $d$ and the stated letters; the states follow from \ref{it:A2}
and \ref{it:P3}; and the breaks at both ends of $P$ survive, because the
site is at least $\beta\ge2d$ away from them. Conversely, let $R'$ be a
piece of $U$ of this class. If a junction, of a piece $P$, lies in $R'$ or
at one of its ends, then $R'$ and $P'$ overlap in at least $d$ positions,
because $R'$ has at least $2d$ positions and $P'$ extends by $2d$ positions
on both sides of its junction; so $R'=P'$ by \ref{it:P1}. Otherwise $R'$ and
its two breaks consist of original text, so $R'$ is the image of a piece of
$W$ of the class, which is not shortened.

A cut $c'$ of $u_l$ is \emph{clean for $\theta$} if no shortened piece $P$ of
$w_l$ with $\cmap_l(s_P)\le c'$ has its junction strictly between $c'$ and
$c'+\theta$. For every $P\in\mathcal C$ and every $\theta$, the cut $\cmap_l(e_P)-d$
is clean for $\theta$, and its greatest preimage is $e_P-d$: a piece $Q'$ that
begins at or before this cut and has its junction after it contains the
last $d$ positions of $P'$, by \ref{it:A1}, so $Q'=P'$ by \ref{it:P1}, but
the junction of $P'$ lies before the cut.

Two words $y,y'$ \emph{agree on $\theta$ positions from the cuts $c,c'$} if
$y[c+x]=y'[c'+x]$ for all $x<\theta$, these positions being defined.

\begin{lemma}[Copying pieces]\label{lem:copy}
Let $u_h$ and $u_g$ agree on $\theta$ positions from the cuts $a'$ and $b'$, and
let $R'$ be a piece of $U$ of the class $\mathcal C$ in $u_h$, between the
cuts $s'$ and $e'$.
\begin{enumerate}
\item\label{it:copy-start} If $a'<s'$ and $s'+2d\le a'+\theta$, then $u_g$ has a
  piece $P'$ of the class that begins at $b'+(s'-a')$, with the same root.
\item\label{it:copy-stretch} If $s'\le a'$ and $a'+\ell\le e'$ for some
  $\ell$ with $2d<\ell\le \theta$, then $u_g$ has a piece $P'$ of the class that
  contains the positions from $b'$ to $b'+\ell$.
\item\label{it:copy-end} In both cases, if $e'<a'+\theta$ then $P'$ ends at
  $b'+(e'-a')$, and otherwise $P'$ contains all positions from its
  beginning to $b'+\theta$.
\end{enumerate}
The same holds with the roles of $u_h$ and $u_g$ exchanged.
\end{lemma}

\begin{proof}
The copy in $u_g$ of the part of $R'$ inside the agreement is a factor of
length at least $2d$ with period $d$ and a primitive conjugate root, so it
lies in a piece of the class. A break of $R'$ whose two letters lie in the
agreement is copied, which determines the corresponding end of $P'$.
\end{proof}

\begin{lemma}[Clean cuts]\label{lem:clean}
Let $T=\wsuf{u_g}{b'}\pre\wsuf{u_h}{a'}$ and $\theta=\len T$. If $b'$ is clean
for $\theta$ in $u_g$ and $a'$ is clean for $\theta$ in $u_h$, then
$\wsuf{w_g}{\gamma(b')}\pre\wsuf{w_h}{\gamma(a')}$.
\end{lemma}

\begin{proof}
By induction on $\theta$. If no junction lies strictly between $a'$ and $a'+\theta$ in
$u_h$ or strictly between $b'$ and $b'+\theta$ in $u_g$, both texts are original,
by \eqref{eq:between-junctions}. Otherwise consider the shortened pieces
whose junctions lie strictly inside one of the two windows. By cleanness
their images begin after $a'$, respectively $b'$; let $r>0$ be the least
distance of such a beginning from the beginning of its window. By
\cref{lem:copy}\eqref{it:copy-start}, applied in one of the two directions,
there are pieces $R'$ in $u_h$ beginning at $a'+r$ and $P'$ in $u_g$
beginning at $b'+r$, with the same root; they are images of pieces $R$ and
$P$ of $\mathcal C$. By the choice of $r$ and by cleanness, no junction lies
strictly between $a'$ and $a'+r$ or at $a'+r$, and likewise in $u_g$; so
\begin{equation}\label{eq:clean-prefix}
  \wsub{w_h}{\gamma(a')}{s_R}=\wsub{u_h}{a'}{a'+r}=\wsub{u_g}{b'}{b'+r}=
  \wsub{w_g}{\gamma(b')}{s_P}.
\end{equation}
By \cref{lem:copy}\eqref{it:copy-end} there are two cases.

\emph{Both pieces end inside the comparison}, at the same distance, so
$L'_P=L'_R$ and $\wsuf{u_g}{\cmap_g(e_P)-d}\pre\wsuf{u_h}{\cmap_h(e_R)-d}$,
a comparison of length less than $\theta$ between clean cuts. By the induction
hypothesis, $\wsuf{w_g}{e_P-d}\pre\wsuf{w_h}{e_R-d}$. In particular the
followers of $P$ and $R$ are comparable. If $P$ and $R$ differed in length
by at least $\kappa$, so would $P'$ and $R'$; so their lengths differ by
less than $\kappa$, $x_P=x_R$ and $L_P=L_R$. Two pieces with the same root
and the same length carry the same text, so
$\wsuf{w_g}{s_P}\pre\wsuf{w_h}{s_R}$.

\emph{$P'$ ends at the end of $u_g$, and $L'_P\le L'_R$.} Then $P$ ends at
the end of $w_g$ and has the empty follower, which is comparable with that
of $R$. If the lengths differ by less than $\kappa$, then $x_P=x_R$;
otherwise the longer piece remains the longer one. In both cases $L_P\le
L_R$, and $\wsuf{w_g}{s_P}$, a word of period $d$ with the root of $R$, is a
prefix of $\wsuf{w_h}{s_R}$.

In both cases the claim follows with \eqref{eq:clean-prefix}.
\end{proof}

\begin{lemma}[Comparisons inside pieces]\label{lem:inside}
Let $P,R\in\mathcal C$ lie in the words $w_g$ and $w_h$, and let
$b'=\cmap_g(s_P)+p_T$ and $a'=\cmap_h(s_R)+p_S$ with $p_T+d\le L'_P$ and
$p_S+d\le L'_R$. Suppose $\wsuf{u_g}{b'}\pre\wsuf{u_h}{a'}$ and
$L'_P-p_T\le L'_R-p_S$, and that $P'$ ends at the end of $u_g$ or
$L'_P-p_T=L'_R-p_S$. Then there are offsets $o_T\le L_P$ and $o_S\le L_R$
such that $\wsuf{w_g}{s_P+o_T}\pre\wsuf{w_h}{s_R+o_S}$ and the states of $W$
at these two cuts are the states of $U$ at $b'$ and $a'$.
\end{lemma}

\begin{proof}
We first show that the followers of $P$ and $R$ are comparable, and that
$\wsuf{w_g}{e_P}\pre\wsuf{w_h}{e_R}$ if $P'$ does not end its word. In that
case $L'_P-p_T=L'_R-p_S$, and shifting the given comparison yields
$\wsuf{u_g}{\cmap_g(e_P)-d}\pre\wsuf{u_h}{\cmap_h(e_R)-d}$ between clean
cuts; \cref{lem:clean} gives $\wsuf{w_g}{e_P-d}\pre\wsuf{w_h}{e_R-d}$. If
$P'$ ends its word, so does $P$, and its follower is empty.

It suffices to find offsets $o_T,o_S$ that are related to $p_T,p_S$ as in
\ref{it:A4} and satisfy $L_P-o_T\le L_R-o_S$, with equality if $P$ does not
end its word. Indeed, the states are then as required; the letters of
$w_g$ from $s_P+o_T$ and of $w_h$ from $s_R+o_S$ agree on $d$ positions,
because they are the letters of $u_g$ from $b'$ and of $u_h$ from $a'$, by
\ref{it:A3} and since $m$ is a multiple of $d$; so by periodicity they
agree up to the end of $P$, and beyond it the followers are compared.

If the lengths of $P$ and $R$ differ by less than $\kappa$, then
$x_P=x_R$, and $o_T=p_T$, $o_S=p_S$ work, since both lengths change by the
same amount. Otherwise the lengths differ by at least $\kappa$, before and
after the deletion and in the same direction, and $\kappa\ge\beta+m\ge
\alpha$.

Suppose $L'_P-p_T=L'_R-p_S$. If $R$ is the longer piece, then
$p_S=p_T+L'_R-L'_P\ge\kappa$; let $o_T=p_T$ and $o_S=p_T+L_R-L_P\ge\kappa$,
which differs from $p_S$ by $m(x_R-x_P)$. If $P$ is the longer piece, let
$o_S=p_S$ and $o_T=p_S+L_P-L_R$, symmetrically.

Suppose $P'$ ends its word and $L'_P-p_T\le L'_R-p_S$. If $P$ is the longer
piece, then $p_T\ge p_S+\kappa\ge\beta$; let $o_T=p_T+m\,x_P$ and
$o_S=p_S$, so that $L_P-o_T=L'_P-p_T\le L'_R-p_S\le L_R-o_S$. If $R$ is the
longer piece, let $o_T=p_T$, and let $o_S=p_S$ if $p_S<\alpha+m$ and
otherwise the number in $[\alpha,\alpha+m)$ that is congruent to $p_S$
modulo $m$. Then $L_R-o_S>L_R-\alpha-m\ge L_P\ge L_P-o_T$.
\end{proof}

\begin{proof}[Proof of \cref{lem:admissible}]
\emph{Domination.} Endpoints and visited states are treated in
\ref{it:A2}. Let $T=\wsuf{u_g}{b'}\pre\wsuf{u_h}{a'}$ be a comparison of
$U$, $\theta=\len T$. We show that the pair of states at $a'$ and $b'$ lies in
$\Chan_{hg}(W)$.

\emph{Case 1: $b'$ is not clean for $\theta$.} Then $b'$ lies in the image $P'$
of a shortened piece, at an offset $p_T<\beta$, and $P'$ continues for more
than $5\beta$ positions after $b'$, inside $T$. By
\cref{lem:copy}\eqref{it:copy-stretch} and \eqref{it:copy-end}, from $u_g$
to $u_h$, $a'$ lies in a piece $R'$ of the class at an offset $p_S$, and
either both pieces end inside the comparison, at the same distance, or $P'$
ends at the end of $u_g$ and $L'_P-p_T\le L'_R-p_S$.
\Cref{lem:inside} applies.

\emph{Case 2: $b'$ and $a'$ are clean for $\theta$.} \Cref{lem:clean} applies,
with the states of \ref{it:A2}.

\emph{Case 3: $b'$ is clean for $\theta$ and $a'$ is not.} Then $a'$ lies in the
image $R'$ of a shortened piece, at an offset $p_S<\beta$, with the junction
of $R'$ strictly inside the comparison. If the comparison continues beyond
the end of $R'$, then by \cref{lem:copy} $b'$ lies in a piece $P'$ of the
class that ends at the same distance, and \cref{lem:inside} applies. So let
$\theta\le L'_R-p_S$. If $\theta>2d$, then by \cref{lem:copy} $T$ lies in a piece $P'$
of the class that ends at the end of $u_g$, and \cref{lem:inside} applies
again. If $\theta\le2d<\beta$, no junction lies strictly inside the window of the
target, since by cleanness it would belong to a piece that begins inside the
window; so $T=\wsuf{w_g}{\gamma(b')}$, and $T$ is read in $w_h$ at the offset
$p_S$ of $R$, with the same state, by \ref{it:A3}.

\emph{First rejected cuts.} Let $W$ be a hole. By \ref{it:A2} and
\cref{lem:hole-dominates}, $U$ is a hole with the same endpoints. The
states of $u_l$ are those of $w_l$ at retained cuts, in the same order, and
the states of $w_l$ are rejected exactly from $J_l$ on; so the first
rejected cut of $u_l$ is $\cmap_l(J_l)$, the image of $J_l$ or of the right
end of the site that contains $J_l$.

\emph{Other classes.} Let $R$ be a piece with a period $d'\le d$, a root
that is not conjugate to $\zeta$, and a length of at least $2d\ge d+d'$. If
a position of a site of $P\in\mathcal C$ lay in $R$ or next to it, then by
\ref{it:A1} and $\beta\ge4d$ the pieces $P$ and $R$ would overlap in at
least $2d$ positions, contradicting \ref{it:P2}. So $R$ and the letters
next to it are not affected, and its image is a piece of $U$ with the same
period, root and length.
\end{proof}

\subsection{The Cap}\label{app:compression-cap}

\begin{lemma}[Simple paths]\label{lem:simple-path}
Let $P_1,\dots,P_v$ be pairwise different pieces of $\mathcal C$ such that
consecutive ones have comparable followers. Suppose that all but at most
$r$ of them belong to a set $\mathcal H$, and that every subset of
$\mathcal H$ with pairwise comparable followers has at most $h$ elements.
Then $v\le(r+1)\,2^h-1$.
\end{lemma}

\begin{proof}
Removing the pieces outside $\mathcal H$ leaves at most $r+1$ paths of
pieces of $\mathcal H$. We show that such a path has at most $2^h-1$
vertices, by induction on $h$. Order its pieces by the prefix order of
their followers, with ties broken by an arbitrary linear order; this is a
partial order in which the elements below any element form a chain, and
consecutive pieces of the path are comparable. So the path has a least
element, which is comparable with all its pieces. Removing it leaves at
most two paths, in each of which pairwise comparable sets have at most
$h-1$ elements, since the least element can be added to them. Hence the
path has at most $1+2(2^{h-1}-1)$ vertices.
\end{proof}

\begin{proof}[Proof of \cref{lem:cap}]
We determine integers $z_P\le0$ and put $x_P=-z_P$. Consider the following
difference constraints, for all pieces $P,Q\in\mathcal C$ with comparable
followers:
\begin{itemize}
\item $z_P\ge z_Q$ and $z_Q\ge z_P$, if the lengths differ by less than
  $\kappa$;
\item $z_Q\ge z_P-\lfloor(L_Q-L_P-\kappa)/m\rfloor$, if $L_P+\kappa\le
  L_Q$;
\end{itemize}
together with the lower bounds $z_P\ge\underline z_P$, where
$\underline z_P=-\lfloor(L_P-\phi)/m\rfloor$ if $L_P\ge\phi$ and
$\underline z_P=0$ otherwise. The zero vector satisfies all of them. A
vector $z\le0$ that satisfies them gives an admissible deletion: a
shortened piece keeps at least $\phi$ letters by the lower bound, and the
second kind of constraint says $L_Q+mz_Q\ge L_P+mz_P+\kappa$.

Let $z$ be the vector whose component $z_P$ is the largest value of
$\underline z_{Q}+c$ over all pieces $Q$ and all paths from $Q$ to $P$ in
the graph of the constraints, $c$ being the weight of the path. Since the
zero vector is a solution, the graph has no cycle of positive weight, so
the maximum exists and is attained by a path without repeated vertices; and
$z$ is a solution with $\underline z\le z\le0$, the least one. Fix $P$ and
such a path $Q=P_1,\dots,P_v=P$. By maximality, $z_{P_1}=\underline
z_{P_1}$ and every constraint along the path is an equality for $z$.

Consider the retained lengths $L'_{P_u}=L_{P_u}+mz_{P_u}$. First,
$L'_{P_1}<\phi+m$. Along a constraint of the first kind the retained length
changes by less than $\kappa$. Along a constraint of the second kind, from
$P_u$ to $P_{u+1}$, equality gives $L'_{P_{u+1}}<L'_{P_u}+\kappa+m$. So
$L'_P<\phi+m+(v-1)(\kappa+m)$. The pieces of the path are pairwise different
and consecutive ones have comparable followers, so $v\le(r+1)2^h-1$ by
\cref{lem:simple-path}, with $\mathcal H$ the set of the pieces that are
not exceptional.
\end{proof}

\subsection{The Passes}\label{app:compression-passes}

A \emph{tail decomposition} of a hole $U$ with the parameters $(\eta,D,E)$
is a set of at most $\eta$ pieces of $U$, the \emph{designated pieces}, with
periods at most $D$ and lengths at least $2D$, such that at most $E$
positions of the rejected tails of the coordinates with immortal endpoints
lie in no designated piece; these are the \emph{uncovered} positions.

\begin{lemma}[One pass]\label{lem:pass}
A hole with a tail decomposition with the parameters $(\eta,D,E)$,
$\eta\ge1$, has a dominating hole with a tail decomposition with the
parameters $(\eta-1,D,E')$, where
\[
  E'=E+\eta\,\Bigl(\phi_D+m_D+\bigl((2\eta+E+1)\,2^{k(N-1)}-1\bigr)
  (\kappa_D+m_D)\Bigr).
\]
\end{lemma}

\begin{proof}
Let $d$ be the largest period of a designated piece, $\zeta$ its root and
$\mathcal C$ the class of $\zeta$. Let $\mathcal H$ consist of the pieces of
$\mathcal C$ that lie in a word with a mortal endpoint or end at or before
the first rejected cut of their word. If two pieces of $\mathcal H$ in one
word, ending at $e<e'$, have comparable followers, then the rest after $e'$
is a prefix of the rest after $e$, so the states at $e$ and $e'$ differ by
\cref{lem:distinct}. Hence a subset of $\mathcal H$ with pairwise
comparable followers has at most $k(N-1)$ elements.

The other pieces of $\mathcal C$ lie in words $u_l$ with immortal endpoints
and end after $J_l$; we claim that there are at most $2\eta+E$ of them. The
last position of such a piece $P$ lies in the tail. Either it is one of the
$E$ uncovered positions, or it lies in a designated piece $R$. In the second
case $P=R$, or $P$ ends in the first $d$ positions of $R$ by \ref{it:P1},
if $R\in\mathcal C$, or $P$ begins before $R$ and ends in the first $2d$
positions of $R$ by \ref{it:P2}, if $R\notin\mathcal C$. Since the ends of
different pieces of $\mathcal C$ in one word are more than $d$ apart, each
designated piece accounts for at most two pieces.

By \cref{lem:cap} there is an admissible deletion after which every piece
of $\mathcal C$ has at most $\phi_d+m_d+((2\eta+E+1)2^{k(N-1)}-1)(\kappa_d+
m_d)$ letters, and this is at most the corresponding value for $D$. By
\cref{lem:admissible} the result $U'$ is a hole that dominates $U$, its
tails consist of the images of the retained positions of the tails of $U$,
and the designated pieces outside $\mathcal C$ are pieces of $U'$ with the
same periods and lengths. We keep these as designated pieces; there are at
most $\eta-1$ of them. A position of a tail of $U'$ that lies in none of them
is the image of an uncovered position of $U$ or lies in the image of one of
the at most $\eta$ designated pieces of $\mathcal C$.
\end{proof}

\begin{proof}[Proof of \cref{thm:tails-bound,thm:tails}]
By \cref{prop:run-cover} a reader with a hole has a hole with a tail
decomposition with the parameters $(g,D,E_0)$, $g=k(\Qc+3)$: a maximal
periodic factor with a period of at most $D$ and a length of at least $2D$
is a piece. Applying \cref{lem:pass} as long as there are designated pieces,
at most $g$ times, gives a hole in which
all positions of the rejected tails of the coordinates with immortal
endpoints are uncovered, and their number is at most $E_g$, because the
sequence $E_r$ is nondecreasing and its recurrence is that of
\cref{lem:pass} with $\eta=g-r$. The closed form is proved in
\cref{app:complexity-closed}, and $\TailB(0,N)=0$ by \cref{sec:tails}.
\end{proof}
